% !TeX root = ../../Book4.tex
% 纲要标题：### 12.1 群同调、群上同调与 bar 消解
% 本轮补充的纲要细目：
% * 群阶可逆时证明 Reynolds 投影、不变子环线性及取不变量的正合性；不把它误作环同态
% * 用模特征增广列的非零一次余圈说明不动元素的提升障碍，并在习题中比较稳定理想的商与多项式模特征例子
% 纲要标签：b4:sec:1101
% 纲要位置：计划纲要.md，第 3496 行。

\section{群同调、群上同调与 bar 消解}
\label{b4:sec:1101}

群的不变量与余不变量分别来自取固定部分和强制认同作用。把它们写成群代数上的 Hom 与张量积，就能用同一份 bar 消解计算各次导出群。

% 纲要内容（仅作写作计划，尚非教材正文）：
%
% * 不变量、余不变量及群代数
% * Ext、Tor 定义与 bar 消解
% * 齐次和非齐次余链
%

\subsection{不变量、余不变量及群代数}
\label{b4:subsec:1101:01}

一个群在模上作用以后，可以只保留被所有群元素固定的向量，也可以把同一轨道上的向量强制认同。这两种做法分别产生子模与商模；前者一般不能保持满射，后者一般不能保持单射。群的上同调与同调记录的正是这两种正合性的缺失。

\begin{definition}\label{b4:def:group-invariants}
设 \(k\) 是交换含幺环，\(G\) 是群。在以 \(G\) 为基的自由 \(k\) 模上，将群乘法作 \(k\) 双线性延拓，所得含幺代数称为 \(G\) 在 \(k\) 上的\term[qundaishu]{群代数}[group algebra]，记为 \(kG\)。其元素是有限和 \(\sum_g a_g g\)，单位元为群单位元所对应的基元。幺性左 \(kG\) 模就是带有 \(k\) 线性 \(G\) 作用的幺性 \(k\) 模。对这样的模 \(M\)，定义
\[
M^G\coloneqq\{\,m\in M\mid gm=m\;\text{对所有}\;g\in G\,\},
\qquad
M_G\coloneqq M/\langle gm-m\mid g\in G,\ m\in M\rangle_k,
\]
分别称为\term[bubianliang]{不变量}[invariants]与\term[yububianliang]{余不变量}[coinvariants]。环同态
\(\varepsilon:kG\to k\)、\(\sum_g a_g g\mapsto\sum_g a_g\) 称为群代数上的\term[zengguang]{增广}[augmentation]；其核记为 \(I_G\)，称为增广理想。通过 \(\varepsilon\) 赋予 \(k\) 平凡的左、右 \(kG\) 模结构。
\end{definition}
\symidx{\(I_G\)：群代数的增广理想}
\symidx{\(M^G,M_G\)：群作用的不变量与余不变量}

\begin{proposition}\label{b4:prop:group-invariant-functors}
设 \(k\) 是交换含幺环，\(G\) 是群，\(M\) 是幺性左 \(kG\) 模，\(k\) 通过增广取平凡左、右 \(kG\) 作用。存在自然同构
\[
\operatorname{Hom}_{kG}(k,M)\cong M^G,
\qquad k\otimes_{kG}M\cong M_G=M/I_GM.
\]
因此不变量函子左正合，余不变量函子右正合。
\end{proposition}
\begin{proof}
映射 \(u:k\to M\) 若为 \(kG\) 线性，则由 \(k\) 线性知 \(u(a)=a u(1)\)，而平凡作用给出
\[
g u(1)=u(g\cdot1)=u(1).
\]
反过来，给定 \(m\in M^G\)，映射 \(u_m(a)=am\) 为 \(k\) 线性，且 \(g u_m(a)=am=u_m(g\cdot a)\)，故为 \(kG\) 线性。因此 \(u\mapsto u(1)\) 与 \(m\mapsto u_m\) 互逆。更一般地，若 \(T\) 是赋予平凡 \(G\) 作用的幺性 \(k\) 模，则 \(k\) 线性等变映射 \(T\to M\) 的像必落在 \(M^G\)，因而唯一经过包含 \(M^G\hookrightarrow M\)。不变量由这种映射的唯一分解刻画。

另一方面，\(I_G\) 作为 \(k\) 模由 \(g-1\) 生成，因为增广为零的元素满足
\[
\sum_g a_g g=\sum_g a_g(g-1).
\]
所以 \(I_GM\) 恰为所有 \(gm-m\) 的 \(k\) 线性生成子模。记 \(q:M\to M_G\) 为商映射，便有 \(q(gm)=q(m)\)。在张量积中，平凡右作用给出同一关系
\[
1\otimes gm=(1\cdot g)\otimes m=1\otimes m.
\]
对一般的 \(r\in kG\)，\(q(rm)=\varepsilon(r)q(m)\)，故双线性映射 \((a,m)\mapsto a q(m)\) 满足 \(kG\) 平衡关系，诱导 \(k\otimes_{kG}M\to M_G\)。上面的张量关系又保证 \(m\mapsto1\otimes m\) 零化 \(I_GM\)，因而下降为反向映射。这两张映射分别为
\[
a\otimes m\longmapsto a\bar m,
\qquad \bar m\longmapsto1\otimes m,
\]
它们互逆，因为 \(1\otimes am=a\otimes m\)。若 \(T\) 如上取平凡作用，\(k\) 线性等变映射 \(M\to T\) 恰好零化所有 \(gm-m\)，所以唯一经过 \(q\)。余不变量由此强制全部群作用成为恒等。

对 \(kG\) 模同态 \(v:M\to N\)，\((vu)(1)=v(u(1))\)，且 \(v(gm-m)=gv(m)-v(m)\)；这些等式使上述同构分别与限制到不变量和下降到商模相容，给出自然性。最后将\cref{b4:prop:hom-left-exact}用于平凡左模 \(k\)，将\cref{b4:prop:tensor-right-exact}用于平凡右模 \(k\)，分别得到两种正合性。
\end{proof}

取 \(k=\mathbb Z\)、\(G=C_2=\langle s\rangle\)，记 \(s\) 作用为取负的整数模为 \(\mathbb Z_{\mathrm{sgn}}\)。固定元素满足 \(m=-m\)，故只有零；而 \(sm-m=-2m\)，故
\[
(\mathbb Z_{\mathrm{sgn}})^G=0,
\qquad (\mathbb Z_{\mathrm{sgn}})_G=\mathbb Z/2\mathbb Z.
\]
给 \(\mathbb F_2\) 平凡作用，约化映射与取负作用相容，所以有短正合列
\[
0\longrightarrow\mathbb Z_{\mathrm{sgn}}\xrightarrow{\times2}
\mathbb Z_{\mathrm{sgn}}\longrightarrow\mathbb F_2\longrightarrow0.
\]
取不变量后得到 \(0\to0\to0\to\mathbb F_2\to0\)，其中第三条箭头 \(0\to\mathbb F_2\) 不满；取余不变量后得到
\[
0\longrightarrow\mathbb Z/2\mathbb Z\xrightarrow{0}\mathbb Z/2\mathbb Z
\xrightarrow{\sim}\mathbb F_2\longrightarrow0,
\]
左端的零映射不单。前一缺失发生在右端，需要右导出；后一缺失发生在左端，需要左导出。

\ProseParagraphBreak
当有限群的阶可以作除数时，上述满射问题有一个直接的解决办法。它既适用于模，也适用于群在多项式环上的作用；后一情形还要注意投影与乘法之间的关系。

\begin{proposition}\label{b4:prop:reynolds-exactness}
设 \(k\) 为交换含幺环，\(G\) 为有限群，且 \(|G|\cdot1_k\) 在 \(k\) 中可逆。对每个幺性左 \(kG\) 模 \(M\)，映射
\[
\mathcal R_M(m)=\frac1{|G|}\sum_{g\in G}gm
\]
是到 \(M^G\) 的自然 \(k\) 线性投影，因而函子 \(M\mapsto M^G\) 正合。若 \(A\) 是带 \(G\) 作用的交换 \(k\) 代数，且每个群元素作用为 \(k\) 代数自同构，则 \(\mathcal R_A\) 还是 \(A^G\) 线性的，即
\[
\mathcal R_A(af)=a\mathcal R_A(f)
\qquad(a\in A^G,\ f\in A).
\]
这个平均投影称为\term[Reynolds-suanzi]{Reynolds 算子}[Reynolds operator]。\footnote{雷诺（Osborne Reynolds，1842-1912），生于贝尔法斯特的英国工程师、物理学家，以流体流动中的雷诺数和湍流研究著称。}
\end{proposition}
\begin{proof}
对 \(h\in G\)，左乘 \(h\) 重排求和项，故 \(h\mathcal R_M(m)=\mathcal R_M(m)\)。若 \(m\in M^G\)，每一项等于 \(m\)，所以 \(\mathcal R_M(m)=m\)。这证明像为 \(M^G\)，且 \(\mathcal R_M^2=\mathcal R_M\)。等变映射与每一项作用交换，因而与平均交换，给出自然性。

由\cref{b4:prop:group-invariant-functors}已知左正合性，只需检查满射。若 \(q:M\to N\) 是满的 \(kG\) 模同态，对 \(n\in N^G\) 取 \(q(m)=n\)，则
\[
q\mathcal R_M(m)=\mathcal R_N(n)=n.
\]
故 \(M^G\to N^G\) 满射。最后，当 \(a\in A^G\) 时，\(g(af)=a\cdot g(f)\)，逐项求和即得 \(A^G\) 线性。
\end{proof}

这种线性不等于保持乘法。例如在特征不为二的域 \(k\) 上，让 \(C_2\) 在 \(k[x]\) 上把 \(x\) 送到 \(-x\)，则 \(\mathcal R(x)=0\)，而 \(\mathcal R(x^2)=x^2\ne\mathcal R(x)^2\)。不变环的乘法来自原代数，不能通过把任意两个元素分别平均后相乘来代替。

\subsection{Ext、Tor 定义与 bar 消解}
\label{b4:subsec:1101:02}

\begin{definition}\label{b4:def:group-cohomology}
设 \(G\) 是群，\(M\) 是幺性左 \(\mathbb ZG\) 模，即带 \(G\) 作用的交换群。对 \(n\ge0\)，定义
\[
H^n(G,M)\coloneqq\operatorname{Ext}_{\mathbb ZG}^n(\mathbb Z,M),
\qquad
H_n(G,M)\coloneqq\operatorname{Tor}^{\mathbb ZG}_n(\mathbb Z,M),
\]
分别称为 \(G\) 以 \(M\) 为系数的\term[qunshangtongdiao]{群上同调}[group cohomology]与\term[quntongdiao]{群同调}[group homology]；Ext 中的 \(\mathbb Z\) 是平凡左模，Tor 中的 \(\mathbb Z\) 是平凡右模。由\cref{b4:prop:group-invariant-functors}及 Ext、Tor 的导出函子定义和平衡性，它们也是
\[
H^n(G,-)=R^n\bigl((-)^G\bigr),
\qquad H_n(G,-)=L_n\bigl((-)_{G}\bigr).
\]
若 \(k\) 是交换含幺环，且系数已给为幺性左 \(kG\) 模，也可用 \(\operatorname{Ext}_{kG}^n(k,M)\)、\(\operatorname{Tor}^{kG}_n(k,M)\) 计算同样的群，其中 \(k\) 分别取平凡左、右作用。
\end{definition}
\symidx{\(H^n(G,M),H_n(G,M)\)：群上同调与群同调}

\cref{b4:def:ext}在第二变量中定义右导出；\cref{b4:thm:tor-balanced}则保证固定平凡右模后，在左模变量中取消解也计算同一个 Tor。因此上面的两条等式分别导出取不变量与取余不变量的函子。周期消解和 bar 余链都是这些导出函子的计算模型。改用 \(kG\) 计算时，不要求 \(M\) 在 \(k\) 上投射：平凡模有一套只用群元素编写的自由消解，在两种底环上给出相同的系数复形。

\begin{proposition}\label{b4:prop:group-bar-resolution}
设 \(k\) 是交换含幺环，\(G\) 是群。对 \(n\ge0\)，令 \(B_n(k,G)\) 为以 \(G^{n+1}\) 为基的自由 \(k\) 模，赋予同时左乘坐标的 \(G\) 作用，得到幺性左 \(kG\) 模。对 \(n\ge1\) 取删除坐标的交错微分 \(\mathrm{d}_n:B_n\to B_{n-1}\)，并取增广 \(\varepsilon:B_0\to k\)：
\[
\mathrm{d}_n(g_0,\ldots,g_n)=\sum_{i=0}^n(-1)^i(g_0,\ldots,\widehat{g_i},\ldots,g_n),
\qquad \varepsilon(g_0)=1.
\]
所得增广复形是平凡幺性左 \(kG\) 模 \(k\) 的自由消解，称为齐次 bar 消解。
\end{proposition}
\begin{proof}
\cref{b4:exm:homogeneous-bar}已证明两次删除成对抵消，并用插入首坐标 \(1\) 的收缩证明底层增广 \(k\) 模复形正合。这里恢复各项的群代数结构。每个元组唯一写成
\[
(g_0,\ldots,g_n)=g_0\cdot(1,g_0^{-1}g_1,\ldots,g_0^{-1}g_n),
\]
所以首坐标为 \(1\) 的元组给出轨道代表，并成为 \(kG\) 基，从而
\[
B_n(k,G)\cong\bigoplus_{(h_1,\ldots,h_n)\in G^n}kG.
\]
收缩是 \(k\) 线性的，一般却不与 \(G\) 作用交换：先左乘 \(g\) 再插入首坐标，所得首项为 \(1\)；先插入再左乘，首项为 \(g\)。这不妨碍检测正合性，因为微分均为 \(kG\) 线性，遗忘作用后其核、像仍是同一底层子模；但不能因此断言增广复形在 \(kG\) 模范畴中可缩。以上自由性与正合性对任意交换含幺环 \(k\) 及任意群 \(G\) 成立，由\cref{b4:thm:ext-balanced}可用它计算 Ext。

若令同一个底层模上的右作用为 \(b\cdot g=g^{-1}b\)，则
\[
(b\cdot g)\cdot h=h^{-1}g^{-1}b=(gh)^{-1}b=b\cdot(gh).
\]
逆元使左作用的次序转换为右作用的次序，直接取 \(b\cdot g=gb\) 一般不满足这个等式。每个左 \(kG\) 基元 \(e\) 所生成的左正则模转成右模后仍自由：同构 \(kG\to kGe\) 将 \(g\) 送到 \(g^{-1}e\)。微分、增广的左线性也保证它们与新右作用相容，故得到平凡右模的自由消解 \(B_\bullet^{\mathrm r}\)。按\cref{b4:def:tor}，群同调由 \(B_\bullet^{\mathrm r}\otimes_{kG}M\) 计算。

对给定的 \(kG\) 系数 \(M\)，这个张量复形与 \(B_\bullet(\mathbb Z,G)^{\mathrm r}\otimes_{\mathbb ZG}M\) 自然同构：两者各次都以同一组右模自由基元配上 \(M\)，微分只使用相同的整数系数与群作用。相应地，\(\operatorname{Hom}_{kG}(B_n(k,G),M)\) 与 \(\operatorname{Hom}_{\mathbb ZG}(B_n(\mathbb Z,G),M)\) 都是由这些自由基元任意取值所得的函数群，预合成微分相同。这就证明了定义末尾的底环说明。
\end{proof}

\begin{proposition}\label{b4:prop:group-homology-degree-zero}
设 \(G\) 是群，\(M\) 是幺性左 \(\mathbb ZG\) 模。则 \(H^0(G,M)=M^G\)、\(H_0(G,M)=M_G\)。每个幺性左 \(\mathbb ZG\) 模的短正合列 \(0\to A\to B\to C\to0\) 都产生群上同调与群同调的长正合列；其中前者的低次部分为
\[
0\longrightarrow A^G\longrightarrow B^G\longrightarrow C^G
\xrightarrow{\partial}H^1(G,A)\longrightarrow H^1(G,B)\longrightarrow\cdots.
\]
\end{proposition}
\begin{proof}
零次结论由\cref{b4:prop:group-invariant-functors}及 Ext、Tor 的零次识别得到。长正合列分别是\cref{b4:thm:ext-long-exact}与\cref{b4:thm:tor-long-exact}在平凡模处的特例。
\end{proof}

\subsection{齐次和非齐次余链}
\label{b4:subsec:1101:03}

Hom 复形中的元素可写成满足
\[
F(tg_0,\ldots,tg_n)=tF(g_0,\ldots,g_n)
\]
的函数 \(F:G^{n+1}\to M\)。这样的函数称为齐次余链；其微分由预合成删除坐标的交错微分得到。为了减少一个自变量，令
\[
[g_1|\cdots|g_n]\coloneqq(1,g_1,g_1g_2,\ldots,g_1\cdots g_n).
\]
这些元素恰好组成 \(B_n\) 的 \(kG\) 基：若首坐标为 \(1\) 的轨道代表是 \((1,h_1,\ldots,h_n)\)，则相邻比值 \(g_1=h_1\)、\(g_i=h_{i-1}^{-1}h_i\) 唯一恢复上面的写法；\(n=0\) 时基元为 \([]=(1)\)。在次数二，直接计算
\[
\begin{aligned}
\mathrm{d}_2[g|h]
&=\mathrm{d}_2(1,g,gh)=(g,gh)-(1,gh)+(1,g)\\
&=g[h]-[gh]+[g],
\end{aligned}
\]
因为 \((g,gh)=g\cdot(1,h)\)。删除首项后产生一个作用系数，删除中间项把 \(g,h\) 合并为 \(gh\)，删除末项则留下 \(g\)。在一般次数中也是这三种变化，所以对 \(n\ge1\)，同一微分写为
\[
\begin{split}
\mathrm{d}_n[g_1|\cdots|g_n]
={}&g_1[g_2|\cdots|g_n]\\
&+\sum_{i=1}^{n-1}(-1)^i[g_1|\cdots|g_i g_{i+1}|\cdots|g_n]
+(-1)^n[g_1|\cdots|g_{n-1}].
\end{split}
\]

\begin{definition}\label{b4:def:group-inhomogeneous-cochains}
设 \(G\) 是群，\(M\) 是幺性左 \(\mathbb ZG\) 模。对 \(n\ge0\)，令 \(C^n(G,M)\) 为全部函数 \(G^n\to M\) 组成的交换群，\(C^0(G,M)=M\)。其中的元素称为\term[feiqici-yulian]{非齐次余链}[inhomogeneous cochain]。定义
\begin{equation}\label{b4:eq:group-cochain-differential}
\begin{split}
(\delta f)(g_1,\ldots,g_{n+1})
={}&g_1f(g_2,\ldots,g_{n+1})\\
&+\sum_{i=1}^{n}(-1)^i f(g_1,\ldots,g_i g_{i+1},\ldots,g_{n+1})\\
&+(-1)^{n+1}f(g_1,\ldots,g_n).
\end{split}
\end{equation}
称 \(Z^n(G,M)=\ker\delta\) 中的元素为余圈，称 \(B^n(G,M)=\operatorname{im}\delta\) 中的元素为余边；次数零的余边规定为零。
\end{definition}
\symidx{\(C^n(G,M)\)：群的非齐次余链群}
\symidx{\(Z^n(G,M),B^n(G,M)\)：群的非齐次余圈与余边}

\begin{proposition}\label{b4:prop:group-cochain-model}
设 \(G\) 是群，\(M\) 是幺性左 \(\mathbb ZG\) 模。齐次与非齐次余链复形在每个次数 \(n\ge0\) 通过下列互逆公式同构：
\[
\begin{split}
f(g_1,\ldots,g_n)&=F(1,g_1,g_1g_2,\ldots,g_1\cdots g_n),\\
F(g_0,\ldots,g_n)&=g_0f(g_0^{-1}g_1,g_1^{-1}g_2,\ldots,g_{n-1}^{-1}g_n).
\end{split}
\]
特别地，\(\delta^2=0\)，且 \(H^n(G,M)=Z^n(G,M)/B^n(G,M)\)。
\end{proposition}
\begin{proof}
给定函数 \(f:G^n\to M\)，在自由基上赋值
\[
\varphi_f([g_1|\cdots|g_n])=f(g_1,\ldots,g_n),
\]
便唯一延拓为 \(\mathbb ZG\) 线性映射 \(\varphi_f:B_n\to M\)。反过来，每个这样的线性映射都由这些取值决定，所以
\[
\operatorname{Hom}_{\mathbb ZG}(B_n,M)\cong C^n(G,M).
\]
一个任意元组先由首坐标 \(g_0\) 移到首坐标为 \(1\) 的代表，再用相邻比值描述该代表，所得取值正是命题的第二式。相邻比值逐次相乘恢复 \(g_0^{-1}g_i\)，故两式互逆；同时左乘 \(t\) 不改变相邻比值，只把首项 \(g_0\) 改为 \(tg_0\)，也直接说明第二式的等变性。

沿用\secref{b4:sec:0802}在\cref{b4:thm:ext-balanced}之前约定的无额外符号的 Hom 微分，第 \(n\) 次微分是
\[
\varphi_f\longmapsto\varphi_f\circ\mathrm{d}_{n+1}.
\]
因此对基元 \([g_1|\cdots|g_{n+1}]\) 求值。bar 微分首项的作用系数 \(g_1\) 可以从群环线性映射中提出，其余各项则直接代入 \(f\)，恰好得到\eqref{b4:eq:group-cochain-differential}。例如在次数一，由刚才的二次 bar 计算得到
\[
\begin{aligned}
(\delta f)(g,h)
&=\varphi_f\bigl(\mathrm{d}_2[g|h]\bigr)\\
&=\varphi_f\bigl(g[h]-[gh]+[g]\bigr)
=g f(h)-f(gh)+f(g).
\end{aligned}
\]
这样两种余链的识别与微分交换。由 \(\mathrm{d}_{n+1}\mathrm{d}_{n+2}=0\) 得 \(\delta^{n+1}\delta^n=0\)，再由 bar 消解与\cref{b4:thm:ext-balanced}计算 Ext，即得上同调结论。
\end{proof}

次数零的计算同样来自 \(\mathrm{d}_1[g]=g[]-[]\)，所以 \(\delta a(g)=ga-a\)，其核恰为不变量。次数一的余圈方程则由上面的 \(\delta f=0\) 得到；其中的 \(g f(h)\) 表明它未必是普通群同态，会受到 \(G\) 在系数上的作用修正。这个修正将解释群扩张中的截面为什么不总能选为同态。

\begin{example}\label{b4:exm:modular-augmentation-obstruction}
设 \(k\) 是特征 \(p>0\) 的域，\(G=C_p\)，\(I_G=\ker(\varepsilon:kG\to k)\)。给 \(k\) 平凡作用，则短正合列
\[
0\longrightarrow I_G\longrightarrow kG\xrightarrow{\varepsilon}k\longrightarrow0
\]
在取不变量后不保持右端满射。事实上，正则作用下的不动元恰为
\(cN\)，其中 \(c\in k\)、\(N=\sum_{g\in G}g\)，因为左乘群元素迫使所有基系数相等。而 \(\varepsilon(cN)=cp=0\)，所以 \(1\in k\) 没有不动提升。

\ProseParagraphBreak
在\cref{b4:prop:group-homology-degree-zero}的连接映射中，把 \(1\) 提升为群代数的单位元 \(1_{kG}\)。由\eqref{b4:eq:group-cochain-differential}，其像由 \(I_G\) 值余圈
\[
c(g)=g-1
\]
表示：\(c(gh)=g\,c(h)+c(g)\)，且每个值的增广均为零。这个类非零。若有 \(a\in I_G\) 使 \(g-1=ga-a\) 对所有 \(g\) 成立，则 \(1-a\) 不动，故存在 \(\lambda\in k\) 使 \(1-a=\lambda N\)。取增广却得到 \(1=\lambda p=0\)，矛盾。因此一次上同调具体记录了这个不动向量不能提升的障碍。
\end{example}
